% ---------------------------------------------------------------------------
% 01 — Estudio de Factibilidad (ES). Compilar desde templates/latex:
%   latexmk -pdf 01-estudio-factibilidad/main-es.tex
% Requiere TeX Live 2024+ (pdflatex + latexmk). Sin rutas absolutas.
% ---------------------------------------------------------------------------
\documentclass[12pt]{article}
\usepackage[spanish,es-tabla]{babel}
\usepackage{../common/preambulo}

% --- Metadatos (una sola fuente de verdad para la portada) ---
\renewcommand{\tituloProyecto}{[Nombre del proyecto]}
\renewcommand{\subtituloProyecto}{Estudio de factibilidad integral}
\renewcommand{\tipoDocumento}{Estudio de factibilidad}
\renewcommand{\autorProyecto}{[Equipo promotor]}
\renewcommand{\fechaProyecto}{[Mes, anio]}
\renewcommand{\versionDoc}{v0.1-borrador}

% --- Parametros financieros nombrados: NUNCA repita cifras a mano en el
% texto (antipatron magic numbers). Cambie el valor aqui y todo el
% documento queda consistente. Cifras de EJEMPLO en MXN — sustituya. ---
\newcommand{\InversionInicial}{4\,000\,000}
\newcommand{\TasaDescuento}{12}
\newcommand{\VANbase}{1\,088\,553}
\newcommand{\TIRbase}{20.95}
\newcommand{\Payback}{3.22}
\newcommand{\VANoptimista}{1\,851\,836}
\newcommand{\VANpesimista}{70\,842}
\newcommand{\TIRpesimista}{12.62}
\newcommand{\Moneda}{MXN}

\begin{document}
% ---------------------------------------------------------------------------
% portada.tex — portada compartida (article). Se incluye con ruta relativa:
%   % ---------------------------------------------------------------------------
% portada.tex — portada compartida (article). Se incluye con ruta relativa:
%   % ---------------------------------------------------------------------------
% portada.tex — portada compartida (article). Se incluye con ruta relativa:
%   \input{../common/portada}
% Requiere los metadatos definidos en el main: \tituloProyecto,
% \subtituloProyecto, \tipoDocumento, \autorProyecto, \fechaProyecto,
% \versionDoc (ver preambulo.sty; todos tienen valor por defecto seguro).
% ---------------------------------------------------------------------------
\begin{titlepage}
  \centering
  \vspace*{2cm}
  {\color{bftCoral}\rule{\textwidth}{2pt}\par}
  \vspace{0.8cm}
  {\large\textsc{\tipoDocumento}\par}
  \vspace{0.6cm}
  {\Huge\bfseries\tituloProyecto\par}
  \vspace{0.5cm}
  {\large\subtituloProyecto\par}
  \vspace{1.2cm}
  {\color{bftCoral}\rule{0.25\textwidth}{2pt}\par}
  \vspace{1.2cm}
  {\large\autorProyecto\par}
  \vspace{0.4cm}
  {\large\fechaProyecto\par}
  \vspace{0.4cm}
  {\small\textsc{\versionDoc}\par}
  \vfill
  {\footnotesize Business Feasibility Toolkit\par}
\end{titlepage}

% Requiere los metadatos definidos en el main: \tituloProyecto,
% \subtituloProyecto, \tipoDocumento, \autorProyecto, \fechaProyecto,
% \versionDoc (ver preambulo.sty; todos tienen valor por defecto seguro).
% ---------------------------------------------------------------------------
\begin{titlepage}
  \centering
  \vspace*{2cm}
  {\color{bftCoral}\rule{\textwidth}{2pt}\par}
  \vspace{0.8cm}
  {\large\textsc{\tipoDocumento}\par}
  \vspace{0.6cm}
  {\Huge\bfseries\tituloProyecto\par}
  \vspace{0.5cm}
  {\large\subtituloProyecto\par}
  \vspace{1.2cm}
  {\color{bftCoral}\rule{0.25\textwidth}{2pt}\par}
  \vspace{1.2cm}
  {\large\autorProyecto\par}
  \vspace{0.4cm}
  {\large\fechaProyecto\par}
  \vspace{0.4cm}
  {\small\textsc{\versionDoc}\par}
  \vfill
  {\footnotesize Business Feasibility Toolkit\par}
\end{titlepage}

% Requiere los metadatos definidos en el main: \tituloProyecto,
% \subtituloProyecto, \tipoDocumento, \autorProyecto, \fechaProyecto,
% \versionDoc (ver preambulo.sty; todos tienen valor por defecto seguro).
% ---------------------------------------------------------------------------
\begin{titlepage}
  \centering
  \vspace*{2cm}
  {\color{bftCoral}\rule{\textwidth}{2pt}\par}
  \vspace{0.8cm}
  {\large\textsc{\tipoDocumento}\par}
  \vspace{0.6cm}
  {\Huge\bfseries\tituloProyecto\par}
  \vspace{0.5cm}
  {\large\subtituloProyecto\par}
  \vspace{1.2cm}
  {\color{bftCoral}\rule{0.25\textwidth}{2pt}\par}
  \vspace{1.2cm}
  {\large\autorProyecto\par}
  \vspace{0.4cm}
  {\large\fechaProyecto\par}
  \vspace{0.4cm}
  {\small\textsc{\versionDoc}\par}
  \vfill
  {\footnotesize Business Feasibility Toolkit\par}
\end{titlepage}

\tableofcontents
\newpage

\section{Resumen ejecutivo y recomendacion}
[2–3 parrafos: que proyecto, cuanto requiere, conclusion Go / No-Go y los
dos indicadores clave: VAN de \$\VANbase{} \Moneda{} y TIR de \TIRbase{}\%.]
La recomendacion se fundamenta en la matriz de decision de la
seccion~\ref{sec:gongo}.

\section{Descripcion del proyecto, alcance y supuestos}
\subsection{Objetivo y alcance}
[Que se pretende lograr; alcance geografico y funcional; que queda
explícitamente FUERA del estudio.]
\subsection{Supuestos criticos}
\begin{itemize}
  \item Supuesto de mercado: [ej. crecimiento anual de la demanda].
  \item Supuesto tecnologico: [ej. disponibilidad de equipos].
  \item Supuesto regulatorio: [ej. permisos considerados].
  \item Supuesto financiero: [tasa de descuento \TasaDescuento{}\%, horizonte 5 anios].
\end{itemize}

\section{Estudio de mercado}
\subsection{Segmentacion y mercado objetivo}
[Perfil de cliente, segmentos, TAM/SAM/SOM estimados.]
\subsection{Oferta y competencia}
[Competidores directos, capacidad instalada, barreras de entrada.]
\subsection{Demanda proyectada}
[Series historicas, encuestas, pronostico a 5 anios.]
\subsection{Analisis FODA}
\begin{table}[h]
  \centering
  \caption{Matriz FODA (resumen)}
  \begin{tabular}{p{0.45\textwidth}p{0.45\textwidth}}
    \toprule
    \textbf{Fortalezas} & \textbf{Debilidades} \\
    \midrule
    {[F1]} & [D1] \\
    \bottomrule
    \textbf{Oportunidades} & \textbf{Amenazas} \\
    \midrule
    {[O1]} & [A1] \\
    \bottomrule
  \end{tabular}
\end{table}

\section{Factibilidad tecnico-operativa}
\subsection{Tecnologia y disenio preliminar}
[Alternativas evaluadas, tecnologia elegida y por que; sitio, capacidad,
infraestructura requerida.]
\subsection{Plan operativo}
[Proceso, mantenimiento, escalamiento, cadena de suministro y proveedores
clave con certificaciones.]

\section{Factibilidad legal y regulatoria (Mexico)}
\begin{table}[h]
  \centering
  \caption{Mapa de obligaciones legales MX (verificar con asesor)}
  \begin{tabular}{p{0.32\textwidth}p{0.32\textwidth}p{0.26\textwidth}}
    \toprule
    \textbf{Ambito} & \textbf{Tramite / norma} & \textbf{Estado} \\
    \midrule
    Constitucion & Acta constitutiva, RFC ante el SAT & [ ] \\
    Operacion & Licencia municipal de funcionamiento, uso de suelo & [ ] \\
    Laboral & Registro patronal IMSS, contratos LFT & [ ] \\
    Ambiental & MIA / LGEEPA si aplica (ver sec.~\ref{sec:amb}) & [ ] \\
    Propiedad industrial & Marcas y patentes ante el IMPI & [ ] \\
    Datos personales & Aviso de privacidad LFPDPPP & [ ] \\
    Sectorial & [Permisos especificos del sector] & [ ] \\
    \bottomrule
  \end{tabular}
\end{table}

\section{Factibilidad organizacional}
[Organigrama propuesto, roles, experiencia del equipo, brechas de
contratacion y plan de cierre de brechas.]

\section{Factibilidad financiera: VAN y TIR}
\subsection{Inversion inicial y flujos proyectados}
Inversion inicial \(I_0 = \$\InversionInicial{}\) \Moneda{}; tasa de
descuento \(r = \TasaDescuento{}\%\); horizonte 5 anios.
\begin{table}[h]
  \centering
  \caption{Flujo de caja proyectado (ejemplo, \Moneda{})}
  \label{tab:flujos}
  \begin{tabular}{crrrrr}
    \toprule
    Concepto & Anio 1 & Anio 2 & Anio 3 & Anio 4 & Anio 5 \\
    \midrule
    Flujo neto & 800\,000 & 1\,200\,000 & 1\,600\,000 & 1\,800\,000 & 2\,000\,000 \\
    Valor presente & 714\,286 & 956\,633 & 1\,138\,848 & 1\,143\,933 & 1\,134\,854 \\
    \bottomrule
  \end{tabular}
\end{table}
\subsection{Indicadores}
\begin{equation}
  \mathrm{VAN} = \sum_{t=1}^{5}\frac{CF_t}{(1+r)^t} - I_0
  = \$\VANbase{}\ \text{\Moneda{}}
  \label{eq:van}
\end{equation}
\begin{equation}
  \mathrm{TIR} = r^* \mid \mathrm{VAN}(r^*) = 0 = \TIRbase{}\%
  \label{eq:tir}
\end{equation}
Periodo de recuperacion (payback simple): \Payback{} anios.
Regla de decision: GO si \(\mathrm{VAN} > 0\) y \(\mathrm{TIR} > r\).
\subsection{Analisis de sensibilidad}
\begin{table}[h]
  \centering
  \caption{Escenarios (flujos optimista +15\%, pesimista --20\%)}
  \begin{tabular}{lrr}
    \toprule
    Escenario & VAN (\Moneda{}) & TIR (\%) \\
    \midrule
    Optimista & \VANoptimista{} & $> \TIRbase{}$ \\
    Base & \VANbase{} & \TIRbase{} \\
    Pesimista & \VANpesimista{} & \TIRpesimista{} \\
    \bottomrule
  \end{tabular}
\end{table}
[Incluso el escenario pesimista mantiene VAN positivo, pero con margen
estrecho: la variable critica es (identificarla).]

\section{Factibilidad ambiental y social}
\label{sec:amb}
[Impactos (emisiones, agua, residuos), plan de mitigacion y monitoreo;
empleo directo/indirecto, aceptacion comunitaria y consulta publica.
Si el proyecto requiere MIA, indicarlo aqui y vincularlo con la tabla
legal.]

\section{Analisis de riesgos}
Marco ISO~31000: identificar, analizar, evaluar y tratar.
\begin{table}[h]
  \centering
  \caption{Registro de riesgos (extracto)}
  \begin{tabular}{p{0.26\textwidth}ccp{0.4\textwidth}}
    \toprule
    Riesgo & Prob. & Impacto & Mitigacion y responsable \\
    \midrule
    {[Retraso de permisos]} & Alta & Critico & [Gestor dedicado; tramite en paralelo] \\
    {[Volatilidad de insumos]} & Media & Alto & [Contratos a precio fijo] \\
    {[Demanda menor a la prevista]} & Media & Alto & [Piloto + contratos preventa] \\
    \bottomrule
  \end{tabular}
\end{table}

\section{Conclusiones y matriz Go / No-Go}
\label{sec:gongo}
\begin{table}[h]
  \centering
  \caption{Matriz de decision ponderada (ejemplo)}
  \begin{tabular}{lccc}
    \toprule
    Criterio & Peso & Puntaje /10 & Aporte \\
    \midrule
    Mercado & 0.25 & [ ] & [ ] \\
    Tecnico-operativo & 0.20 & [ ] & [ ] \\
    Legal MX & 0.10 & [ ] & [ ] \\
    Organizacional & 0.10 & [ ] & [ ] \\
    Financiero (VAN/TIR) & 0.25 & [ ] & [ ] \\
    Ambiental-social y riesgos & 0.10 & [ ] & [ ] \\
    \midrule
    \textbf{Total} & 1.00 & & \textbf{[ ]} \\
    \bottomrule
  \end{tabular}
\end{table}
Regla: GO si total $\geq 7.0$ y ningun criterio $< 5.0$.
\begin{bftGoBox}
  Declarar aqui la decision final con condiciones (ej. sujeta a mitigar
  los riesgos R1 y R2 y a confirmar el supuesto S1).
\end{bftGoBox}

\end{document}
